\section{Relativistic stellar structure}
\label{sec:stars}

\subsection{Background configurations}

For a static spherical star,
\begin{equation}
    ds^2
    =
    -e^{2\nu(r)}dt^2
    +
    \left(1-\frac{2m(r)}{r}\right)^{-1}dr^2
    +
    r^2d\Omega^2 .
\end{equation}
The Tolman--Oppenheimer--Volkoff equations are
\begin{align}
    \frac{dm}{dr}
    &=4\pi r^2\epsilon,
    \\
    \frac{dp}{dr}
    &=
    -\frac{(\epsilon+p)(m+4\pi r^3p)}
    {r(r-2m)} .
    \label{eq:tov}
\end{align}
We use relativistic enthalpy
\begin{equation}
    h(p)=\int_0^p\frac{dp'}{\epsilon(p')+p'}
    \label{eq:enthalpy}
\end{equation}
as the integration coordinate. This fixes the stellar surface at \(h=0\) and makes the domain convenient for differentiating stellar observables with respect to EOS degrees of freedom.

The surface values define
\begin{equation}
    M=m(R),
    \qquad
    C=M/R .
\end{equation}

\subsection{Tidal deformability and moment of inertia}

The quadrupolar tidal deformability is
\begin{equation}
    \lambda_2=\frac{2}{3}k_2R^5,
    \qquad
    \Lamtwo=\frac{\lambda_2}{M^5}
    =\frac{2}{3}k_2C^{-5}.
    \label{eq:lambda}
\end{equation}
We integrate the standard even-parity \(l=2\) static perturbation equation together with the background and evaluate \(k_2\) from the surface logarithmic derivative and compactness. \todo{Insert the exact enthalpy-coordinate perturbation equation once the implementation is frozen so the manuscript cannot drift from the production code.}

For uniform rotation at first order in angular velocity, the Hartle frame-dragging equation determines the angular momentum \(J\), from which
\begin{equation}
    I=\frac{J}{\Omega},
    \qquad
    \Ibar=\frac{I}{M^3}.
\end{equation}
The exact enthalpy-coordinate system used for the production calculation will likewise be inserted after the benchmark stage.

\subsection{EOS coordinates and perturbations}

We require EOS perturbations that approach a functional description while satisfying thermodynamic stability and causality by construction. The baseline representation uses a latent field \(u(h)\) mapped to the sound speed,
\begin{equation}
    c_s^2(h)=\sigma[u(h)],
    \qquad
    \sigma(u)=\frac{1}{1+e^{-u}},
    \label{eq:cs-map}
\end{equation}
with stricter bounds applied explicitly if required. The latent field is expanded as
\begin{equation}
    u(h)=u_0(h)+\sum_{i=1}^{N_b}a_i\phi_i(h).
    \label{eq:basis}
\end{equation}
The basis \(\phi_i\) is a numerical representation only. Physical conclusions must converge as \(N_b\) and the basis family are varied.

Given \(c_s^2=dp/d\epsilon\), the remaining thermodynamic functions are reconstructed consistently from an anchor point and joined to an adopted low-density crust treatment. \todo{Freeze the exact thermodynamic reconstruction, density domain, and crust prescription after Stage 1 validation.}

For the EOS response metric we will use at least two choices for \(C_{\rm EOS}\): a smooth covariance in the latent field and a second covariance with shorter correlation length or broader localized freedom. A response hierarchy that exists only for one EOS metric will not be interpreted as universal.
